\documentclass[reprint,amsmath,amssymb,aps]{revtex4-2}
\usepackage{graphicx}
\usepackage{bm}
\usepackage{hyperref}

\begin{document}

\title{Comprehensive Analysis of 4th- and 6th-Order Summation-By-Parts Operators \\ in Spherical Coordinates}

\author{Antigravity Computational Framework}
\affiliation{Advanced Agentic Coding}
\date{\today}

\begin{abstract}
This scientific report details the theoretical foundation, methodological construction, and programmatic implementation of 4th-order and 6th-order Summation-By-Parts (SBP) finite difference operators in Julia (\texttt{sbp4.jl} and \texttt{sbp6.jl}). These numerical operators strictly mimic the continuous integration-by-parts principle on discretely bounded domains, ensuring discrete energy stability for initial-boundary value problems (IBVPs) through dynamically solved rational linear constraint systems.
\end{abstract}

\maketitle

\section{Introduction and Theoretical Context}

The Summation-By-Parts (SBP) framework is a discrete, systematic mechanism for constructing highly stable numerical schemes for partial differential equations (PDEs), such as initial-boundary value problems (IBVPs). By mirroring strict continuous functional properties analytically at the discrete stage, numerical stability flows naturally via the global energy method without resorting to ad-hoc, destabilizing numerical dissipation parameters.

\subsection{The Continuous Integration-By-Parts Principle}
In the continuous functional space, the generic 1D Integration-By-Parts (IBP) property for bounded, continuously differentiable functions $u(x)$ and $v(x)$ on a radial spatial domain $[0, R]$ is algebraically stated as:
\begin{equation}
    \int_0^R u \frac{\partial v}{\partial x} dx = \left[ u \cdot v \right]_0^R - \int_0^R \frac{\partial u}{\partial x} v dx
\end{equation}

\subsection{The Discrete SBP Form}
An underlying finite difference first-derivative operator $D$ satisfies the strict SBP property if it can be analytically factored into the universal form:
\begin{equation}
    D = P^{-1}Q
\end{equation}
where:
\begin{itemize}
    \item $P$ is a symmetric positive-definite (SPD) matrix defining the discrete volume norm (quadrature). The inner product is $(u, v)_P = u^T P v \approx \int u v dx$.
    \item $Q$ is an ``almost'' skew-symmetric matrix such that its symmetric part extracts the exact physical boundary evaluations:
\end{itemize}
\begin{equation}
    Q + Q^T = B = \text{diag}(-1, 0, \dots, 0, 1)
\end{equation}

Evaluating the discrete derivative through the inner product yields $(u, Dv)_P + (Du, v)_P = u^T B v$, structurally confirming discrete energy conservation identical to the analytical IBP boundary condition model.

\subsection{Diagonal vs. Block-Diagonal Norms and Boundary Accuracy}
Traditionally, a perfectly diagonal norm $P$ requires the interior difference stencil yielding $2s$-order accuracy to drop steeply to $s$-order accuracy across boundary closures. If one upgrades $P$ to an extended block-diagonal structure (where nodes strictly near the boundary interact non-diagonally), higher accuracy boundary closures can be preserved (commonly $2s-1$). 

The implementations analyzed herein (\texttt{sbp4.jl} and \texttt{sbp6.jl}) utilize an advanced \textit{split-pattern formulation} primarily tailored toward spherical/radial coordinates incorporating $r^p$ metric terms. Instead of a single unmanageable non-diagonal $P$, the system separates into a strictly diagonal scalar mass matrix $S$ evaluated alongside a structured, split vector mass matrix $V$. The localized divergence becomes:
\begin{equation}
    D = S^{-1}(B - G_{\text{even}}^T V)
\end{equation}
This architecture retains peak memory sparsity everywhere except within the explicitly specified coupling subset within $V$ localized tightly near the origin boundary $r=0$.

\section{4th-Order SBP Operator (\texttt{sbp4.jl})}

The module \texttt{sbp4.jl} establishes the 4th-order operator utilizing a tightly banded vector-mass modification block near the physical boundaries.

\subsection{Interior Stencil}
For the vast majority of the interior grid (nodes strictly unaffected by the boundary perturbation matrices), the localized divergence exactly replicates the standardized 4th-order Cartesian central difference stencil. Evaluated at a node $x_i$ with grid spacing $h$, the finite-difference operator $G_{\text{full}}$ acts as:
\begin{equation}
    (G_{\text{full}} u)_i = \frac{1}{12h} \left( u_{i-2} - 8u_{i-1} + 8u_{i+1} - u_{i+2} \right) + \mathcal{O}(h^4)
\end{equation}

\subsection{Boundary Closure and Variable Point Modifications}
Because $P$ must be non-trivially manipulated to satisfy the formal energy bounds, boundary nodes $i \in [1, 5]$ are logically relaxed into free numerical algebraic symbols. 

In \texttt{sbp4.jl}, the scalar mass $S$ remains diagonal but holds 5 free parameters for boundary modification: $s_1, \dots, s_5$. For indices $i \ge 6$, the scalar masses identically match the unperturbed metric transformation $s_i = (H_{\text{cart\_half}} r^p)_i$.

The vector mass matrix $V$ carries a localized symmetric tridiagonal-style block near the origin for indices up to 5:
\begin{itemize}
    \item Fixed Origin: $V_{1,1} = 1$
    \item Diagonal Unknowns: $V_{2,2}, V_{3,3}, V_{4,4}, V_{5,5}$
    \item Off-Diagonal Unknown Pairs: $V_{2,3}=V_{3,2}$, $V_{3,4}=V_{4,3}$, $V_{4,5}=V_{5,4}$
\end{itemize}
For indices greater than 5, the diagonal of $V$ directly conforms to $S$, meaning $V_{i,i} = s_i \ \forall i \ge 6$, and strictly isolates boundary cross-coupling.

\subsection{Programmatic Operator Imposition}
Rather than hardcoding arbitrary floating-point boundary coefficients, the Julia implementation executes an exact rational linear solver (\texttt{sbp4\_solve\_accuracy\_constraints}) solving $A \mathbf{x} = \mathbf{b}$ populated by constraint polynomials:
\begin{enumerate}
    \item \textbf{Uniform Linear Gradients:} $D \cdot r = (p+1)$ mapping uniform slopes across $N-4$ nodes.
    \item \textbf{Polynomial Derivative Moments:} $D \cdot r^3 = (p+3)r^2$ imposed physically over the first 5 rows (\texttt{rows\_r3 = collect(1:min(5, N))}).
    \item \textbf{Global Mass Integral Accuracy:} $\mathbf{1}^T S \mathbf{1} = \frac{R^{p+1}}{p+1}$ tracking quadrature sum invariants.
\end{enumerate}


\section{6th-Order SBP Operator (\texttt{sbp6.jl})}

The secondary module \texttt{sbp6.jl} natively incorporates the substantially broader stencil sampling geometry requisite to force global 6th-order error scaling, $\mathcal{O}(h^6)$, dictating an exceptionally expansive boundary perturbation layer.

\subsection{Interior Stencil Configuration}
The interior folded spatial gradient $G_{\text{even}}$ maps precisely to the 6th-order interior central difference derivative projected into rational matrices:
\begin{equation}
    (G_{\text{full}} u)_i = \frac{ -u_{i-3} + 9u_{i-2} - 45u_{i-1} + 45u_{i+1} - 9u_{i+2} + u_{i+3} }{60h}
\end{equation}

\subsection{Expanded Boundary Closure Framework}
To critically counterbalance the mathematically dominant $u_{i \pm 3}$ nodes propagating beyond the coordinate origin boundaries, the degrees of algebraic freedom inside the $V$ boundary sub-matrix are heavily expanded (as defined via \texttt{sbp6\_v\_offdiag\_pairs(N)}).
\begin{itemize}
    \item Diagonal perturbation parameters now range through $V_{2,2}$ to \textbf{$V_{10,10}$}.
    \item The localized boundary block transitions functionally to a \textbf{pentadiagonal bandwidth pattern}. For nodes $i \in [2,5]$, interactions natively cross-track radial coupling distances of 1 and 2: $(i, i+1)$ and $(i, i+2)$. Examples include $(2,3)$ and $(2,4)$, up through $(5,7)$. Row 6 supports only an edge-capping pair $(6,7)$.
\end{itemize}

\subsection{Constraints and Augmented Methodology}
Given the vastly extended stencil limits, \texttt{sbp6.jl} introduces significant programmatic complexity regarding constraint mapping to retain positive-definiteness natively. To construct stable bounds mapped across 6th-order global interpolations, the script aggressively extends structural polynomial invariants:
\begin{verbatim}
add_degree_constraints!(5, rows.rows_r5, :divergence)
\end{verbatim}
This systematically imposes $D \cdot r^5 = (p+5)r^4$ against the first $m$ origin nodes. By enforcing highly erratic, non-linear algebraic derivative limits, the code generates stable pentadiagonal boundaries retaining exact mathematical precision under finite limits. Moreover, all $N$ scalar diagonal masses within $S$ are temporarily freed within the constraint domain during the sparse factorization process to preserve matrix compliance limits.


\section{Algorithmic Methodology and Optimization}

\subsection{Matrix Formulation with Exact Relational Arithmetic}
The central optimized innovation within these modules lies specifically in circumventing generic IEEE-754 64-bit floating point approximations. Classical inverted boundaries theoretically drift from exact symmetry by factors exceeding $10^{-16}$, degrading eigenvalue limits over extended high-frequency time integration iterations.
\begin{verbatim}
Rq = _sbp4_as_big_rational(R)
\end{verbatim}
Both \texttt{sbp4.jl} and \texttt{sbp6.jl} natively promote all geometrical factors, constraints, and difference geometries into absolute \texttt{Rational\{BigInt\}} components.

\subsection{Parity Folding and Symmetrical Base Gradients}
Instead of deriving incredibly complex, numerically ill-conditioned one-sided differences at boundary nodes, the framework implements an elegant Parity Folding scheme. The code evaluates exact symmetrical operations scaling from $-R$ to $R$, capturing boundary parity via analytical parity matrices $E_{\text{even}}$ (for even base functions) and $E_{\text{odd}}$ (for odd nodes).
\begin{verbatim}
Geven = sparse(Rop * Gfull * Eeven)
\end{verbatim}
This functionally projects origin conditions out to symmetrical analytical geometry, guaranteeing boundary symmetry inherently rather than iteratively.

\subsection{Matrix-Based Optimization Assembly}
All algebraic manipulations resolve perfectly analytically via linear solvers against parameter subsets, and uniquely, are populated \textit{only subsequently} into normalized Sparse arrays utilizing standard library constructors:
\begin{verbatim}
S = spdiagm(0 => sdiag)
I, J, GV = findnz(sparse(Geven))
Gq = sparse(...)
\end{verbatim}
This separates arbitrary theoretical scaling logic against high-performance runtime evaluations, rendering matrix-free scaling entirely sparse and bounded, ideal for rapid computational PDE integrators running large radial models structurally dependent on stable properties.

\end{document}
